\documentclass[11pt,a4paper]{article}
\usepackage{amsmath}
\usepackage{amsfonts}
\usepackage{amssymb}
\usepackage{graphicx}
\usepackage{float}

\title{Local Randomized Neural Networks with Hybridized Discontinuous Petrov-Galerkin Methods for Stokes-Darcy Flows}
\author{Summary of Dang and Wang (2023)}
\date{}

\begin{document}

\maketitle

\section{Introduction}

The Stokes-Darcy equations describe the coupling of fluid flow in a porous medium with a free flow region, which has significant applications in environmental science, biomedical engineering, and industrial processes. The coupling is governed by the Beavers-Joseph law or its simplified version, the Beavers-Joseph-Saffman law, which specifies interface conditions between the two regions.

Solving these coupled flow problems presents two main challenges: (1) designing numerical schemes that are stable for both Stokes-dominated and Darcy-dominated regions, and (2) handling the interface conditions to effectively couple different numerical methods across domains. Traditional approaches using finite element methods, discontinuous Galerkin methods, and weak Galerkin methods have been applied with varying degrees of success.

This paper introduces a novel approach combining local randomized neural networks (LRNNs) with the hybridized discontinuous Petrov-Galerkin (HDPG) method to solve both Stokes-Darcy problems and Brinkman equations. The method partitions the domain into subdomains, constructs an LRNN on each subdomain, and uses the HDPG scheme to couple these neural networks to approximate the unknown functions.

\section{Mathematical Representation of Randomized Neural Networks}

\subsection{Fully Connected Feedforward Neural Networks}

A fully connected feedforward neural network applies linear transformations and nonlinear activation functions sequentially to input data. The network consists of $L+2$ layers, where $L$ is the number of hidden layers. Each layer $l$ has $m_l$ neurons, with weights $W^{(l)} \in \mathbb{R}^{m_l \times m_{l-1}}$ and biases $b^{(l)} \in \mathbb{R}^{m_l \times 1}$ connecting the $(l-1)$-th and $l$-th layers.

Mathematically, the network is expressed as:
\begin{align}
y^{(0)} &= x, \\
y^{(l)} &= \rho(W^{(l)}y^{(l-1)} + b^{(l)}), \quad l = 1, \ldots, L, \\
z &= W^{(L+1)}y^{(L)} + b^{(L+1)},
\end{align}
where $x \in \mathbb{R}^{m_0}$ is the input, $y^{(l)}$ are the hidden layers, and $z$ is the output.

\subsection{Randomized Neural Networks}

Randomized neural networks (RNNs) differ by randomly assigning and fixing the values of hidden layer parameters during training. Only the output layer weights are trained, typically by solving a least-squares problem. The hidden layer parameters include weights $W = \{W^{(1)}, \ldots, W^{(L)}\}$ and biases $B = \{b^{(1)}, \ldots, b^{(L)}\}$.

The output of the last hidden layer defines a set of basis functions $\phi^{W,B}_\rho(x) = y^{(L)}$. A randomized neural network function space can then be defined as:
\begin{equation}
[N^\theta_\rho(D)]^n = \left\{z(x) = W\phi^{W,B}_\rho(x) : x \in D \subset \mathbb{R}^d, \forall W \in \mathbb{R}^{n \times N}\right\},
\end{equation}
where $\theta = \{W, B\}$ represents the fixed parameters, $d$ is the input dimension, $N$ is the number of neurons in the last hidden layer, and $n$ is the output dimension.

\section{LRNN-HDPG Method for Darcy Flows}

\subsection{Problem Formulation}

The Darcy problem with Dirichlet boundary conditions is given by:
\begin{align}
\mathbf{u} &= -K\nabla p \quad \text{in } \Omega, \\
\nabla \cdot \mathbf{u} &= f \quad \text{in } \Omega, \\
p &= g \quad \text{on } \partial\Omega,
\end{align}
where $\mathbf{u}$ is the fluid velocity, $p$ is the pressure, $K$ is the permeability-to-viscosity ratio, $f$ is the source term, and $g$ is the boundary condition.

\subsection{Hybridized Discontinuous Galerkin Formulation}

The domain $\Omega$ is partitioned into elements $\mathcal{T}_h$, with edges/faces denoted by $\mathcal{E}_h$. Function spaces are defined:
\begin{align}
V &= \left\{\mathbf{v} \in [L^2(\Omega)]^d : \mathbf{v}|_K \in [L^2(K)]^d, \forall K \in \mathcal{T}_h \right\}, \\
\hat{V} &= \left\{\hat{\mathbf{v}} \in [L^2(\mathcal{E}_h)]^d : \hat{\mathbf{v}} = \hat{v}\mathbf{n}_e, \hat{v}|_e \in L^2(e), \forall e \in \mathcal{E}_h \right\}, \\
Q &= \left\{q \in L^2(\Omega) : q|_K \in H^1(K), \forall K \in \mathcal{T}_h \right\}.
\end{align}

The hybridized DG formulation seeks $\mathbf{u} \in V$, $\hat{\mathbf{u}} \in \hat{V}$, and $p \in Q$ such that:
\begin{align}
(K^{-1}\mathbf{u}, \mathbf{v})_{\mathcal{T}_h} + (\nabla p, \mathbf{v})_{\mathcal{T}_h} - (p, \hat{\mathbf{v}} \cdot \mathbf{n})_{\partial\mathcal{T}_h} &= -(g, \hat{\mathbf{v}} \cdot \mathbf{n})_{\mathcal{E}_h^b}, \\
(\mathbf{u}, \nabla q)_{\mathcal{T}_h} - (\hat{\mathbf{u}} \cdot \mathbf{n}, q)_{\partial\mathcal{T}_h} &= -(f, q)_{\mathcal{T}_h},
\end{align}
for all $(\mathbf{v}, \hat{\mathbf{v}}) \in V \times \hat{V}$ and $q \in Q$.

\subsection{LRNN-HDPG Method}

The LRNN-HDPG method discretizes the hybridized DG formulation using RNN spaces for trial functions and piecewise polynomial spaces for test functions.

Trial function spaces:
\begin{align}
V_r &= \left\{\mathbf{u}_r \in [L^2(\Omega)]^d : \mathbf{u}_r|_K \in [N^\theta_\rho(K)]^d, \forall K \in \mathcal{T}_h \right\}, \\
\hat{V}_r &= \left\{\hat{\mathbf{u}}_r \in [L^2(\mathcal{E}_h)]^d : \hat{\mathbf{u}}_r = \hat{u}_r\mathbf{n}_e, \hat{u}_r|_e \in N^\theta_\rho(e), \forall e \in \mathcal{E}_h \right\}, \\
Q_r &= \left\{p_r \in L^2(\Omega) : p_r|_K \in N^\theta_\rho(K), \forall K \in \mathcal{T}_h \right\}.
\end{align}

Test function spaces:
\begin{align}
V_h &= \left\{\mathbf{v}_h \in [L^2(\Omega)]^d : \mathbf{v}_h|_K \in [P_k(K)]^d, \forall K \in \mathcal{T}_h \right\}, \\
\hat{V}_h &= \left\{\hat{\mathbf{v}}_h \in [L^2(\mathcal{E}_h)]^d : \hat{\mathbf{v}}_h = \hat{v}_h\mathbf{n}_e, \hat{v}_h|_e \in P_k(e), \forall e \in \mathcal{E}_h \right\}, \\
Q_h &= \left\{q_h \in L^2(\Omega) : q_h|_K \in P_{k+1}(K), \forall K \in \mathcal{T}_h \right\}.
\end{align}

The LRNN-HDPG scheme with a stabilization term seeks $\mathbf{u}_r \in V_r$, $\hat{\mathbf{u}}_r \in \hat{V}_r$, and $p_r \in Q_r$ such that:
\begin{align}
(K^{-1}\mathbf{u}_r, \mathbf{v}_h)_{\mathcal{T}_h} + \eta_h((\mathbf{u}_r - \hat{\mathbf{u}}_r) \cdot \mathbf{n}, \mathbf{v}_h \cdot \mathbf{n})_{\partial\mathcal{T}_h} + (\nabla p_r, \mathbf{v}_h)_{\mathcal{T}_h} &= 0, \\
-\eta_h((\mathbf{u}_r - \hat{\mathbf{u}}_r) \cdot \mathbf{n}, \hat{\mathbf{v}}_h \cdot \mathbf{n})_{\partial\mathcal{T}_h} - (p_r, \hat{\mathbf{v}}_h \cdot \mathbf{n})_{\partial\mathcal{T}_h} &= -(g, \hat{\mathbf{v}}_h \cdot \mathbf{n})_{\mathcal{E}_h^b}, \\
(\mathbf{u}_r, \nabla q_h)_{\mathcal{T}_h} - (\hat{\mathbf{u}}_r \cdot \mathbf{n}, q_h)_{\partial\mathcal{T}_h} &= -(f, q_h)_{\mathcal{T}_h},
\end{align}
for all $\mathbf{v}_h \in V_h$, $\hat{\mathbf{v}}_h \in \hat{V}_h$, and $q_h \in Q_h$, where $\eta \geq 0$ is a stabilization parameter.

Two key variations of the method are also presented:

1. **Degrees of Freedom Reduction**: The velocity $\mathbf{u}_r$ can be eliminated from the system using an operator defined by the first equation, reducing the computational cost.

2. **Global Trace Neural Network**: Instead of defining $\hat{\mathbf{u}}_r$ piecewise on edges/faces, a single neural network can be used across the entire domain and evaluated only at edges/faces.

\section{LRNN-HDPG Method for Stokes-Darcy Flows}

\subsection{Stokes Equations in Velocity-Stress Formulation}

The Stokes equations with velocity-stress formulation are:
\begin{align}
\sigma^d - 2\nu\varepsilon(\mathbf{u}) &= 0 \quad \text{in } \Omega, \\
-\nabla \cdot \sigma &= \mathbf{f} \quad \text{in } \Omega, \\
\mathbf{u} &= \mathbf{g} \quad \text{on } \partial\Omega, \\
\int_\Omega \text{tr}(\sigma) dx &= 0,
\end{align}
where $\sigma^d = \sigma - \frac{1}{d}\text{tr}(\sigma)I$ is the deviatoric stress, $\varepsilon(\mathbf{u}) = (\nabla \mathbf{u} + (\nabla \mathbf{u})^t)/2$ is the strain rate, $\nu$ is the viscosity, and pressure $p$ can be recovered as $p = -\frac{1}{d}\text{tr}(\sigma)$.

The hybridized DG formulation is discretized similarly to the Darcy case, leading to the LRNN-HDPG scheme for Stokes equations.

\subsection{Coupled Stokes-Darcy Problem}

For the coupled problem, the domain is divided into $\Omega_S$ (Stokes region) and $\Omega_D$ (Darcy region) with an interface $\Gamma$. The interface conditions are:
\begin{align}
\mathbf{u}^S \cdot \mathbf{n} &= \mathbf{u}^D \cdot \mathbf{n} \quad \text{on } \Gamma, \\
(\sigma^S\mathbf{n}) \cdot \mathbf{n} &= -p^D \quad \text{on } \Gamma, \\
(\sigma^S\mathbf{n}) \cdot \mathbf{t} &= -\nu\kappa^{-1}(\mathbf{u}^S \cdot \mathbf{t}) \quad \text{on } \Gamma,
\end{align}
where $\mathbf{n}$ and $\mathbf{t}$ are unit normal and tangent vectors, and $\kappa$ is the friction coefficient.

A key innovation of the LRNN-HDPG method is handling these interface conditions by directly enforcing them at sampling points on $\Gamma$. The discrete versions at sampling points $\mathbf{x}_i \in \Gamma$ are:
\begin{align}
\mathbf{u}^S_r(\mathbf{x}_i) \cdot \mathbf{n} &= \hat{\mathbf{u}}^D_r(\mathbf{x}_i) \cdot \mathbf{n}, \\
(\hat{\sigma}^S_r(\mathbf{x}_i)\mathbf{n}) \cdot \mathbf{n} &= -p^D_r(\mathbf{x}_i), \\
(\hat{\sigma}^S_r(\mathbf{x}_i)\mathbf{n}) \cdot \mathbf{t} &= -\nu\kappa^{-1}(\mathbf{x}_i)(\mathbf{u}^S_r(\mathbf{x}_i) \cdot \mathbf{t}).
\end{align}

The complete system combines the LRNN-HDPG schemes for Stokes and Darcy regions with these interface conditions.

\subsection{Brinkman Equations}

The Brinkman equations provide another approach to coupling Stokes and Darcy flows:
\begin{align}
\nu\kappa^{-1}\mathbf{u} - \nabla \cdot (2\nu\varepsilon(\mathbf{u})) + \nabla p &= \mathbf{f} \quad \text{in } \Omega, \\
\nabla \cdot \mathbf{u} &= 0 \quad \text{in } \Omega, \\
\mathbf{u} &= \mathbf{g} \quad \text{on } \partial\Omega,
\end{align}
where $\kappa$ is the permeability tensor. The flow is dominated by Darcy's law when $\kappa$ approaches zero and by Stokes flow when $\kappa$ is large.

The LRNN-HDPG method for the Brinkman equations follows a similar formulation as for the Stokes equations with an additional term accounting for the Darcy resistance.

\section{Numerical Examples and Results}

The numerical implementation used RNNs with a single hidden layer and the hyperbolic tangent activation function. All numerical results were averaged over multiple experiments to account for randomness in the neural network parameters.

\subsection{Darcy Flows Through Anisotropic Porous Media}

The first example considered the Darcy problem with anisotropic permeability tensor $K = \begin{bmatrix} 10 & 2 \\ 2 & 100 \end{bmatrix}$ and exact solution $p(x,y) = x(1-x)y(1-y)e^{xy}$.

Results showed:
\begin{itemize}
    \item The LRNN-HDPG method achieved better accuracy with increased degrees of freedom and higher-order test functions.
    \item The stabilization term did not enhance results, so subsequent experiments used $\eta=0$.
    \item The degree of freedom reduction scheme maintained accuracy while reducing computational cost.
    \item The global trace neural network approach yielded accuracy comparable to the standard approach.
    \item Compared to the weak Galerkin mixed finite element method, the LRNN-HDPG method achieved more accurate solutions with fewer degrees of freedom.
\end{itemize}

\subsection{Darcy Problem with Highly Oscillatory Solutions}

This example tested the method on problems with highly oscillatory solutions:
$p(x,y) = \frac{1}{m}\sum_{i=1}^{m}\sin(2i\pi x)\sin(2i\pi y)$, where $m$ is a positive integer.

Results demonstrated:
\begin{itemize}
    \item The LRNN-HDG method (using RNNs for both trial and test functions) improved accuracy with more degrees of freedom and finer meshes.
    \item For increasing oscillation complexity (larger $m$), the method required finer meshes but still achieved high accuracy.
    \item Compared to physics-informed neural networks based on multilevel domain decomposition, the LRNN-HDG method achieved more accurate solutions.
\end{itemize}

\subsection{Stokes Equations with Different Viscosities}

This example tested the method on the Stokes equations with viscosities $\nu = 0.1, 0.01$, and $0.001$.

Results showed:
\begin{itemize}
    \item The LRNN-HDPG method achieved good results across all tested viscosity values.
    \item Higher degrees of freedom and finer meshes consistently improved accuracy.
    \item Relative $L^2$ errors as low as $10^{-5}$ were achieved for all variables (stress, velocity, and pressure).
\end{itemize}

\subsection{Stokes-Darcy Equations with Beavers-Joseph Law}

This example evaluated the coupled Stokes-Darcy problem with the Beavers-Joseph interface condition, testing various combinations of permeability coefficient $\alpha$ (0.01 and 100) and viscosity $\nu$ (0.001 and 0.1).

Results showed:
\begin{itemize}
    \item The LRNN-HDPG method performed well for all parameter combinations.
    \item Interface conditions were effectively enforced through sampling points.
    \item The method handled both the simplified Beavers-Joseph-Saffman law and the more general Beavers-Joseph law.
\end{itemize}

\subsection{Brinkman Equations}

The final example tested the Brinkman equations with varying permeability-inverse coefficient $\alpha = 10^{-3}, 1$, and $10^3$.

Results demonstrated:
\begin{itemize}
    \item The LRNN-HDPG method performed well across all regimes (Stokes-dominated to Darcy-dominated).
    \item Higher polynomial degrees and finer meshes improved accuracy.
    \item The method achieved relative $L^2$ errors below $10^{-4}$ for all variables in the optimal cases.
\end{itemize}

\section{Implications and Advantages}

\subsection{Computational Efficiency}

\begin{itemize}
    \item The LRNN-HDPG method reduces training complexity by fixing hidden layer parameters randomly and only training output layer weights through least-squares methods, avoiding computational burden and convergence issues associated with nonconvex optimization.
    \item The method achieves high accuracy with relatively few degrees of freedom compared to traditional methods.
    \item The degrees of freedom reduction scheme offers further computational savings.
\end{itemize}

\subsection{Stability and Flexibility}

\begin{itemize}
    \item The method is stable for both Stokes-dominated and Darcy-dominated regimes without requiring different numerical schemes for different domains.
    \item Interface conditions are handled elegantly through direct enforcement at sampling points, avoiding complex coupling strategies.
    \item The method accommodates both the Beavers-Joseph-Saffman law and the more general Beavers-Joseph law without modification.
    \item Different activation functions can be used for different elements, enhancing adaptability to complex problems.
\end{itemize}

\subsection{Unified Framework}

\begin{itemize}
    \item The LRNN-HDPG method provides a unified approach for solving Darcy flows, Stokes flows, coupled Stokes-Darcy problems, and Brinkman equations.
    \item The neural network structure can be adjusted to handle different types of PDEs while maintaining the same overall framework.
    \item The method naturally extends to other coupled problems in computational mechanics.
\end{itemize}

\subsection{Future Research Directions}

\begin{itemize}
    \item Extension to nonlinear partial differential equations
    \item Improved neural network structures to enhance approximation capabilities
    \item Parallel processing implementation for enhanced efficiency
    \item Mesh adaptation strategies for complex problems
    \item Development of rigorous error estimates for the methods
\end{itemize}

\end{document} 